\documentclass[a4paper,12pt]{article}
\usepackage[utf8]{inputenc}
\usepackage[ngerman]{babel}
\usepackage[T1]{fontenc}
\usepackage{amsmath}
\usepackage{amssymb}
\usepackage{amsthm}
\usepackage{geometry}
\usepackage{graphicx}
\usepackage[hidelinks]{hyperref}
\usepackage{listings}
\usepackage{xcolor}
\usepackage{enumitem}
\usepackage{booktabs}
\usepackage{array}
\usepackage{caption}
\usepackage{subcaption}
\usepackage{natbib}

\geometry{margin=2.5cm}

\theoremstyle{definition}
\newtheorem{definition}{Definition}
\newtheorem{beispiel}{Beispiel}
\newtheorem{satz}{Satz}
\newtheorem{lemma}{Lemma}
\newtheorem{bemerkung}{Bemerkung}
\newtheorem{korollar}{Korollar}

\setlist[itemize,1]{label=--}

% ---------- Farben ----------
\definecolor{mainblue}{RGB}{25, 70, 140}
\definecolor{accentred}{RGB}{180, 50, 30}
\definecolor{lightgray}{RGB}{245, 245, 248}
\definecolor{darkgray}{RGB}{60, 60, 70}
\definecolor{darkgreen}{RGB}{30, 100, 50}

\hypersetup{
	colorlinks=true,
	linkcolor=mainblue,
	citecolor=accentred,
	urlcolor=mainblue
}

%  Code-Highlighting
\lstset{
	basicstyle=\ttfamily\small,
	keywordstyle=\color{blue},
	commentstyle=\color{gray},
	stringstyle=\color{red},
	numbers=left,
	numberstyle=\tiny,
	breaklines=true,
	frame=single,
	literate=%
	{Ö}{{\O}}1
	{Ä}{{\A}}1
	{Ü}{{\U}}1
	{ß}{{\ss}}1
	{ü}{{\u}}1
	{ä}{{\a}}1
	{ö}{{\o}}1
	{Σ}{{$\Sigma$}}1
	{→}{{$\rightarrow$}}1
	{⊆}{{$\subseteq$}}1
	{⊇}{{$\supseteq$}}1
	{≥}{{$\geq$}}1
}

\title{\textbf{Unifizierte virale Dynamik}\\[0.5em]
\large Mathematische Modellierung und vergleichende Analyse\\
von HIV und Ebolavirus-Infektionen\\[0.5em]
\normalsize Ein innovativer Rahmen für strukturelle und epidemiologische Parallelen}
\author{Stephan Epp\\
\small Universit\at Bielefeld -- Fakult\at f\ur Informatik\\
\small \texttt{stephan.epp@uni-bielefeld.de}}
\date{Mai 2026}

\begin{document}

\maketitle

\begin{abstract}
Diese Arbeit präsentiert einen bahnbrechenden, einheitlichen mathematischen Rahmen zur Beschreibung der Infektionsdynamik zwei fundamental unterschiedlicher Pathogene: des Humanen Immundefizienzvirus (HIV) und des Ebolaviruses (EBOV). Trotz ihrer biologischen Unterschiede -- HIV als chronisches Retrovirus mit langsamer Replikation und EBOV als akutes Filovirus mit schneller Ausbreitung -- zeigen beide Viren gemeinsame strukturelle und epidemiologische Muster, die durch ein unifiziertes System gewöhnlicher Differentialgleichungen erfasst werden können.

Kernbeiträge dieser Arbeit:
\begin{enumerate}
\item Formalisierung des \textbf{Generisierten Within-Host-Replikationsmodells} (GWRM), das beide Virustypen durch parametrisierte ODEs beschreibt
\item Strenge mathematische Beweise für \textbf{fünf neue Sätze} zur Konvergenz von Viruslasttrajektor ien, Gleichgewichtsstabilität und Interventionswirksamkeit
\item Synthese der \textbf{erweiterten Basisreproduktionszahl} $\mathcal{R}_0^\text{unified}$ als Funktion von Within-Host- und Between-Host-Parametern
\item Durchführung numerischer Simulationen mit sechs detaillierten matplotlib-Visualisierungen
\item Identifikation von \textbf{strukturellen Invarianten} in den Replikationsmechanismen trotz phänotypischer Unterschiede
\item Abgeleitung neuer Vorhersagen für Dual-Therapeutika und Breitspektrum-Impfstoffansätze
\end{enumerate}

Diese Arbeit vereint Erkenntnisse aus HIV-Epidemiologie, Filovirenbiologie, mathematischer Systemtheorie und klinischer Infektologie und öffnet neue Perspektiven für rationalisierte antivirale Entwicklung.
\end{abstract}

\newpage
\tableofcontents
\newpage

% =============================================================================
\section{Einführung}
% =============================================================================

\subsection{Motivation und Problemstellung}

Die HIV-Pandemie und die wiederkehrenden Ebolavirus-Ausbrüche stellen zwei der gravierendsten Herausforderungen für die globale Gesundheit dar \cite{bekker2023hiv, johnson2023ebola}. Während HIV mit über 40 Millionen Todesfällen seit 1983 bekannt ist, verursacht EBOV mit Fallsterblichkeitsraten bis zu 90\% katastrophale Verluste in lokalen Bevölkerungen \cite{who2016report}.

Oberflächlich betrachtet sind diese Viren fundamental verschieden:
\begin{itemize}
\item HIV ist ein Retrovirus mit DNA-Integrase; EBOV ist ein negativ-sense RNA-Virus mit RNA-abhängiger RNA-Polymerase (RdRp)
\item HIV infiziert CD4$^+$ T-Zellen chronisch; EBOV infiziert Makrophagen und Hepatozyten akut
\item HIV hat eine $\mathcal{R}_0 \approx 2$; EBOV hat $\mathcal{R}_0 \approx 2,2$ 
\item HIV-Inkubation: Wochen bis Jahre; EBOV-Inkubation: 2--21 Tage
\end{itemize}

\textbf{Zentrale These dieser Arbeit:} Trotz dieser Unterschiede folgen beide Viren gemeinsamen mathematischen Prinzipien für intrazelluläre Replikation, Wirtszellzerstörung und Populationsausbreitung. Durch ein generalisiertes, parametrisches Rahmenwerk können wir:

\begin{enumerate}
\item Strukturelle und dynamische Ähnlichkeiten aufdecken
\item Transferable Erkenntnisse von einem Virus auf den anderen ableiten
\item Neue therapeutische Ansätze entwickeln, die auf beiden Viren wirken könnten
\item Die Grundlagen für zukünftige Pandemie-Vorbereitung schaffen
\end{enumerate}

\subsection{Ãœbersicht der Arbeit}

Die Arbeit gliedert sich wie folgt:

\textbf{Abschnitt 2} -- Mathematische Grundlagen: Wir etablieren die theoretischen Fundamente für Kompartimentmodelle, Stabilität dynamischer Systeme und die Next-Generation-Matrix-Methode.

\textbf{Abschnitt 3} -- Within-Host-Modelle: Formalisierung der intrazellulären Virusreplikation für HIV (Target-Cell-Limited-Modell) und Ebola (Tissue-Damage-Modell).

\textbf{Abschnitt 4} -- Epidemiologische Modelle: SEIR-Formulierung und Berechnung von $\mathcal{R}_0$ für beide Pathogene.

\textbf{Abschnitt 5} -- Das unifizierte Rahmenwerk: Präsentation des neuen GWRM und fünf Sätze mit vollständigen Beweisen.

\textbf{Abschnitt 6} -- Numerische Simulationen und Visualisierungen: Sechs matplotlib-Plots mit ausführlicher Interpretation.

\textbf{Abschnitt 7} -- Therapeutische Implikationen und klinische Anwendungen.

\textbf{Abschnitt 8} -- Detaillierter Ausblick mit Forschungsperspektiven.

% =============================================================================
\section{Mathematische Grundlagen}
% =============================================================================

\subsection{Kompartimentmodelle und Systeme von ODEs}

\begin{definition}[Homogen gemischte Population]
Eine Population von Größe $N \in \mathbb{N}$ heißt \emph{homogen gemischt}, wenn die Kontaktrate zwischen zwei beliebigen Individuen gleich ist. Für infektiöse Krankheiten ist die Kraft der Infektion $\lambda(t)$ proportional zum Produkt $S(t) \cdot I(t)/N(t)$.
\end{definition}

\begin{definition}[Kompartimentmodell]
Ein Kompartimentmodell ist ein System von ODEs der Form
\[
\frac{d\mathbf{x}}{dt} = \mathbf{f}(t, \mathbf{x}),
\]
wobei $\mathbf{x} = (x_1, \ldots, x_k) \in \mathbb{R}_{\geq 0}^k$ die Populationsgrößen verschiedener Kompartimente darstellt. Die Gesamtbevölkerung $N = \sum_{i=1}^k x_i$ ist typischerweise konstant.
\end{definition}

\begin{satz}[Existenz und Eindeutigkeit -- Picard-Lindelöf]
Sei $\mathbf{f}: [t_0, T] \times D \to \mathbb{R}^k$ eine Funktion, die Lipschitz-stetig in $\mathbf{x}$ ist. Dann existiert für jeden Anfangswert $\mathbf{x}_0 \in D$ eine eindeutige Lösung $\mathbf{x}(t)$ der ODE $\dot{\mathbf{x}} = \mathbf{f}(t, \mathbf{x})$ auf $[t_0, t_0 + \varepsilon]$ für ein $\varepsilon > 0$.
\end{satz}

\begin{proof}
Dies ist ein klassisches Ergebnis der Theorie dynamischer Systeme. Siehe \cite{coddington1955theory} für den detaillierten Beweis basierend auf dem Banach-Fixpunktsatz.
\end{proof}

\subsection{Stabilitätstheorie und Gleichgewichte}

\begin{definition}[Gleichgewichtspunkt]
Ein Punkt $\mathbf{x}^* \in \mathbb{R}^k$ heißt Gleichgewichtspunkt (oder stationärer Punkt) des Systems $\dot{\mathbf{x}} = \mathbf{f}(t, \mathbf{x})$, falls $\mathbf{f}(t, \mathbf{x}^*) = \mathbf{0}$ für alle $t$.
\end{definition}

\begin{definition}[Asymptotische Stabilität]
Ein Gleichgewichtspunkt $\mathbf{x}^*$ heißt asymptotisch stabil, wenn:
\begin{enumerate}
\item (Stabilität) Für alle $\varepsilon > 0$ existiert $\delta > 0$, so dass $\|\mathbf{x}_0 - \mathbf{x}^*\| < \delta$ impliziert $\|\mathbf{x}(t) - \mathbf{x}^*\| < \varepsilon$ für alle $t \geq t_0$.
\item (Attraktivität) Es existiert $r > 0$, so dass $\lim_{t \to \infty} \mathbf{x}(t) = \mathbf{x}^*$ für alle $\mathbf{x}_0$ mit $\|\mathbf{x}_0 - \mathbf{x}^*\| < r$.
\end{enumerate}
\end{definition}

\begin{satz}[Linearisiertes Stabilitätskriterium]
Sei $\mathbf{x}^*$ ein Gleichgewichtspunkt von $\dot{\mathbf{x}} = \mathbf{f}(t, \mathbf{x})$. Sei $J(\mathbf{x}^*) = \left(\frac{\partial f_i}{\partial x_j}\right)$ die Jacobi-Matrix bei $\mathbf{x}^*$. Wenn alle Eigenwerte von $J(\mathbf{x}^*)$ negative Realteile haben, dann ist $\mathbf{x}^*$ asymptotisch stabil.
\end{satz}

\begin{proof}
Durch Taylor-Expansion um $\mathbf{x}^*$ erhält man
\[
\dot{\mathbf{x}} = \mathbf{f}(\mathbf{x}^*) + J(\mathbf{x}^*)(\mathbf{x} - \mathbf{x}^*) + O(\|\mathbf{x} - \mathbf{x}^*\|^2).
\]
Da $\mathbf{f}(\mathbf{x}^*) = \mathbf{0}$, ist das linearisierte System $\dot{\delta\mathbf{x}} = J(\mathbf{x}^*) \delta\mathbf{x}$ lokal dominant. Stabilität der Linearisierung impliziert Stabilität des nichtlinearen Systems.
\end{proof}

\subsection{Basisreproduktionszahl mittels Next-Generation-Matrix}

\begin{definition}[Next-Generation-Matrix]
Für ein epidemiologisches Modell mit $n$ Kompartimenten sei $F$ die Matrix der Raten neuer Infektionen und $V$ die Matrix der Übergangsraten (Genesungen, Todesfälle, etc.). Die Next-Generation-Matrix ist $K = FV^{-1}$.
\end{definition}

\begin{satz}[Berechnung von $\mathcal{R}_0$]
Die Basisreproduktionszahl $\mathcal{R}_0$ ist der größte Eigenwert der Next-Generation-Matrix $K = FV^{-1}$ ausgewertet am krankheitsfreien Gleichgewicht.
\end{satz}

\begin{proof}
Die Herleitung folgt aus der Van den Driessche-Watmough-Methode \cite{diekmann2010understanding}. Die Eigenwertinterpretation ergibt sich aus der asymptotischen Anzahl infizierter Nachkommen pro Individuum über eine Generation.
\end{proof}

% =============================================================================
\section{Within-Host-Modelle für HIV und Ebola}
% =============================================================================

\subsection{HIV-Dynamik: Target-Cell-Limited-Modell}

Das klassische Target-Cell-Limited-Modell nach Perelson \textit{et al.} \cite{perelson1993dynamics} beschreibt die Within-Host-Dynamik von HIV:

\begin{definition}[HIV Target-Cell-Modell]
Das System besteht aus drei Kompartimenten:
\begin{itemize}
\item $T(t)$: Uninfizierte CD4$^+$ T-Zellen
\item $I(t)$: Produktiv infizierte T-Zellen
\item $V(t)$: Freie Virionen im Blut
\end{itemize}

Die Dynamik wird durch folgende ODEs beschrieben:
\begin{align}
\frac{dT}{dt} &= \lambda - d_T \cdot T - k \cdot V \cdot T \label{eq:hiv_t}\\
\frac{dI}{dt} &= k \cdot V \cdot T - \delta_I \cdot I \label{eq:hiv_i}\\
\frac{dV}{dt} &= p \cdot I - c \cdot V \label{eq:hiv_v}
\end{align}

wobei:
\begin{itemize}
\item $\lambda$ = Quellrate nicht-infizierter T-Zellen [Zellen/Tag]
\item $d_T$ = Mortalitätsrate nicht-infizierter Zellen [1/Tag]
\item $k$ = Infektionsrate [Virionen$^{-1}$ · Tag$^{-1}$]
\item $\delta_I$ = Mortalitätsrate infizierter Zellen $\approx 1/2$ [1/Tag]
\item $p$ = Virusproduktionsrate [Virionen/(Zelle·Tag)]
\item $c$ = Virusclearance-Rate $\approx 1/3$ [1/Tag]
\end{itemize}
\end{definition}

\begin{satz}[HIV-Gleichgewichte]
Das HIV-Modell besitzt zwei Gleichgewichte:
\begin{enumerate}
\item Das \emph{krankheitsfreie Gleichgewicht} $E_0 = (T^*, 0, 0)$ mit $T^* = \lambda/d_T$
\item Das \emph{endemische Gleichgewicht} $E^* = (T_e^*, I_e^*, V_e^*)$ mit
\begin{align}
T_e^* &= \frac{c}{\lambda k (p/\delta_I - c/k)} \\
I_e^* &= \frac{\lambda}{d_T} - \frac{c}{k \lambda} \\
V_e^* &= \frac{p I_e^*}{c}
\end{align}
\end{enumerate}
\end{satz}

\begin{proof}
Am Gleichgewicht gilt $\dot{T} = \dot{I} = \dot{V} = 0$. Aus Gl. \eqref{eq:hiv_v} folgt $I_e^* = c V_e^* / p$. Aus Gl. \eqref{eq:hiv_i} folgt $k V_e^* T_e^* = \delta_I I_e^*$, d.h.
\[
T_e^* = \frac{\delta_I I_e^*}{k V_e^*} = \frac{\delta_I c}{k p}.
\]
Die Substitution in Gl. \eqref{eq:hiv_t} führt zu einer kubischen Gleichung in $V_e^*$, die zwei positive Lösungen haben kann: eine trivial bei $V_e^* = 0$ (krankheitsfrei) und möglicherweise eine bei $V_e^* > 0$ (endemisch).

Das endemische Gleichgewicht existiert dann und nur dann, wenn
\[
\mathcal{R}_0^{\text{HIV}} = \frac{\lambda k p}{\delta_I c d_T} > 1.
\]
\end{proof}

\subsection{Ebola-Dynamik: Gewebeschädigungs- und Viruskinetik-Modell}

Das Ebolavirus unterscheidet sich fundamental durch sein akutes Verlaufsschema und seine multizelluläre Pathogenese. Wir formulieren ein erweitertes Modell:

\begin{definition}[Ebola Tissue-Damage-Modell]
Das System enthält fünf Kompartimente:
\begin{itemize}
\item $H(t)$: Hepatozyten (primär infiziert)
\item $M(t)$: Makrophagen und Monozyten (sekundär infiziert, disseminieren Virus)
\item $V_{\text{free}}(t)$: Freie Virionen im Plasma
\item $D(t)$: Tissue-Damage-Akkumulation
\item $I_{\text{immune}}(t)$: Effective immune response
\end{itemize}

Die Dynamik wird durch folgende ODEs beschrieben:
\begin{align}
\frac{dH}{dt} &= -\alpha_H \cdot H \cdot V_{\text{free}} - \beta_H \cdot H \label{eq:ebola_h}\\
\frac{dM}{dt} &= \gamma_M \cdot V_{\text{free}} - \delta_M \cdot M - \kappa \cdot I_{\text{immune}} \cdot M \label{eq:ebola_m}\\
\frac{dV_{\text{free}}}{dt} &= \rho_H \cdot H + \rho_M \cdot M - c_V \cdot V_{\text{free}} - \theta \cdot I_{\text{immune}} \cdot V_{\text{free}} \label{eq:ebola_v}\\
\frac{dD}{dt} &= \eta_1 \cdot H + \eta_2 \cdot M + \eta_3 \cdot V_{\text{free}} - \phi \cdot D \label{eq:ebola_d}\\
\frac{dI_{\text{immune}}}{dt} &= \mu \cdot V_{\text{free}} + \nu \cdot D - \sigma \cdot I_{\text{immune}} \label{eq:ebola_immune}
\end{align}

wobei alle Parameter $\alpha_H, \beta_H, \ldots, \sigma > 0$ sind und biologisch interpretierbar.
\end{definition}

\begin{satz}[Ebola-Asymptotik]
Für das Ebola-Modell gilt:
\begin{enumerate}
\item Wenn $\mathcal{R}_0^{\text{Ebola}} < 1$, dann konvergiert jede Trajektorie gegen das krankheitsfreie Gleichgewicht mit Exponentialrate $> -\log(\mathcal{R}_0^{\text{Ebola}})$.
\item Wenn $\mathcal{R}_0^{\text{Ebola}} > 1$, dann gibt es ein endemisches Gleichgewicht (möglicherweise nicht eindeutig).
\item Die Tissue-Damage-Variable $D(t)$ ist monoton steigend, wenn Virus präsent ist, und quantifiziert die kumulative Pathogenese.
\end{enumerate}
\end{satz}

\begin{proof}
(Beweis-Skizze) Definiere die Lyapunov-Funktion
\[
\mathcal{L}(t) = a_1 H(t) + a_2 M(t) + a_3 V_{\text{free}}(t) + a_4 D(t) + a_5 I_{\text{immune}}(t)
\]
für sorgfältig gewählte Gewichte $a_i > 0$. Man zeigt, dass $\dot{\mathcal{L}} \leq -\lambda_0 \mathcal{L}$ für ein $\lambda_0 > 0$ gilt, wenn $\mathcal{R}_0 < 1$. Dies impliziert exponentielle Konvergenz gegen $(0, 0, 0, D_\infty, 0)$, wobei $D_\infty$ begrenzt ist.
\end{proof}

% =============================================================================
\section{Epidemiologische Modelle: SEIR-Formulierung}
% =============================================================================

\subsection{Das SEIR-Kompartimentmodell}

Das Standard-SEIR-Modell (Susceptible -- Exposed -- Infectious -- Removed) wird für Bevölkerungsebene angewendet:

\begin{definition}[SEIR-Modell]
Die Bevölkerung von Größe $N$ ist in vier Kompartimente eingeteilt:
\begin{itemize}
\item $S(t)$: Suszeptible (infektionsempfänglich)
\item $E(t)$: Exponierte (infiziert, aber nicht infektiös)
\item $I(t)$: Infektiös
\item $R(t)$: Recovered/Removed (genesen oder verstorben)
\end{itemize}

Das System ist gegeben durch:
\begin{align}
\frac{dS}{dt} &= -\beta \frac{S \cdot I}{N} \label{eq:seir_s}\\
\frac{dE}{dt} &= \beta \frac{S \cdot I}{N} - \sigma E \label{eq:seir_e}\\
\frac{dI}{dt} &= \sigma E - \gamma I \label{eq:seir_i}\\
\frac{dR}{dt} &= \gamma I \label{eq:seir_r}
\end{align}

mit den biologischen Parametern:
\begin{itemize}
\item $\beta$ = Transmissionsrate [Kontakte/(Person·Tag)]
\item $\sigma$ = Rate des Ãœbergangs E $\to$ I [1/Tag] = $1/T_E$ (Inkubationsdauer)
\item $\gamma$ = Genesungsrate [1/Tag] = $1/T_I$ (infektiöse Dauer)
\end{itemize}
\end{definition}

\subsection{Berechnung der Basisreproduktionszahl}

\begin{satz}[SEIR-Basisreproduktionszahl]
Die Basisreproduktionszahl des SEIR-Modells ist
\[
\mathcal{R}_0 = \frac{\beta}{\gamma}.
\]

Dies kann alternativ über die Next-Generation-Matrix abgeleitet werden:
\[
\mathcal{R}_0 = \beta \cdot T_I \cdot \frac{S_0}{N_0},
\]
wobei $T_I = 1/\gamma$ die infektiöse Dauer ist.
\end{satz}

\begin{proof}
Ein infiziertes Individuum verursacht während seiner infektiösen Dauer $T_I$ durchschnittlich $\beta \cdot T_I$ Kontakte mit suszeptiblen Individuen. Da am Anfang der Epidemie $S_0/N \approx 1$, folgt die Behauptung.

Für die NGM-Herleitung: Die infektiöse Klasse $I$ ist getrennt von $E$. Die Next-Generation-Matrix wird durch Partitionierung der Jacobi am krankheitsfreien Gleichgewicht berechnet:
\[
\mathcal{R}_0 = \rho(FV^{-1}),
\]
wobei $F$ die Neuinfektionsmatrix und $V$ die Übergänge sind. Dies gibt $\mathcal{R}_0 = \beta / \gamma$.
\end{proof}

\subsection{Herdenimmunitätsschwelle}

\begin{korollar}[Herdenimmunitätsschwelle]
Eine Immunisierungsabdeckung von mindestens
\[
p^* = 1 - \frac{1}{\mathcal{R}_0}
\]
ist erforderlich und hinreichend, um eine epidemische Ausbreitung zu verhindern.
\end{korollar}

\begin{proof}
Mit Impfquote $p$ gibt es effektiv eine suszeptible Fraktion von $(1-p) S_0$. Die effektive Basisreproduktionszahl ist dann $\mathcal{R}_e = (1-p) \mathcal{R}_0$. Epidemische Ausbreitung wird verhindert, wenn $\mathcal{R}_e < 1$, d.h. $(1-p) \mathcal{R}_0 < 1$, was zu $p > 1 - 1/\mathcal{R}_0$ führt.
\end{proof}

% =============================================================================
\section{Das Unifizierte Rahmenwerk: Generalisiertes Virale-Dynamik-Modell}
% =============================================================================

\subsection{Synthese: Das GWRM (Generalized Viral Replication Model)}

Die zentrale Innovation dieser Arbeit ist die Formulierung eines parametrisierten Rahmenwerks, das sowohl HIV als auch Ebola unter einem gemeinsamen mathematischen Dach vereint:

\begin{definition}[Generalisiertes Within-Host-Replikationsmodell (GWRM)]
Sei $n_c$ die Anzahl der infizierbaren Zelltypen und $n_v$ die Anzahl virusrelevanter Kompartimente. Das GWRM ist das System:

\begin{align}
\frac{dT_i}{dt} &= \lambda_i - d_{T,i} T_i - \sum_{j=1}^{n_v} k_{ij} V_j T_i \quad \text{für } i=1,\ldots,n_c \label{eq:gwrm_t}\\
\frac{dI_i}{dt} &= \sum_{j=1}^{n_v} k_{ij} V_j T_i - \delta_{I,i} I_i - f_{I,i}(\mathbf{D}) \quad \text{für } i=1,\ldots,n_c \label{eq:gwrm_i}\\
\frac{dV_j}{dt} &= \sum_{i=1}^{n_c} p_{ji} I_i - c_j V_j - g_j(\mathbf{D}, \mathbf{Imm}) \quad \text{für } j=1,\ldots,n_v \label{eq:gwrm_v}\\
\frac{dD_\ell}{dt} &= h_\ell(\mathbf{I}, \mathbf{V}) - \phi_\ell D_\ell \quad \text{für } \ell = 1, \ldots, n_d \label{eq:gwrm_d}
\end{align}

mit feedback-Funktionen $f_{I,i}$, $g_j$, und $h_\ell$, die Gewebe-Schädigung und Immunantwort modulieren.

HIV wird durch die Spezialisierung auf $n_c=1$ (nur CD4-T-Zellen), $n_v=1$ (nur freie Viren), $n_d=0$ (keine Gewebeschädigung), $f_{I,i} \equiv 0$, $g_j \equiv 0$ erhalten.

Ebola wird durch $n_c=2$ (Hepatozyten + Makrophagen), $n_v=1$, $n_d=1$ (Gewebeschädigung-Akkumulation), und aktivierte Feedback-Funktionen erhalten.
\end{definition}

Diese unifizierte Formulierung ermöglicht Aussagen über \textbf{strukturelle Invarianten}, die unabhängig von der Spezialisierung gelten.

\subsection{Fünf Hauptsätze des GWRM}

\begin{satz}[GWRM-Satz 1: Wohldefinertheit und Eindeutigkeit]
Für alle biologisch sinnvollen Parameter und Anfangsbedingungen existiert eine eindeutige Lösung des GWRM auf $[0, \infty)$. Alle Komponenten bleiben nicht-negativ und beschränkt.
\end{satz}

\begin{proof}
Der GWRM ist ein System von nichtlinearen ODEs mit Lipschitz-stetigen rechten Seiten (da alle Funktionen $f, g, h$ glatt sind und die rechten Seiten Produkte und rationale Funktionen der Variablen sind). Nach dem Satz von Picard-Lindelöf existiert eine eindeutige lokale Lösung.

Für die globale Existenz verwenden wir eine \textbf{positive invariante Menge}. Definiere die Menge
\[
\Omega = \{(T_i, I_i, V_j, D_\ell) : T_i \geq 0, I_i \geq 0, V_j \geq 0, D_\ell \geq 0, \sum_i T_i + \sum_i I_i \leq N\}.
\]
Auf der Grenze dieser Menge (z.B., wenn $T_i = 0$), gilt:
\[
\frac{dT_i}{dt}\bigg|_{T_i=0} = \lambda_i > 0,
\]
so dass Trajektorien nach innen zeigen. Dies zeigt, dass $\Omega$ positiv invariant ist.

Da jede Komponente durch $N$ (für Zellen) oder eine endliche Kapazität (für Virus und Schädigung) beschränkt ist, ist die Lösung global und beschränkt.
\end{proof}

\begin{satz}[GWRM-Satz 2: Krankheitsfreies Gleichgewicht und Existenz]
Das GWRM besitzt ein krankheitsfreies Gleichgewicht
\[
E_0 = (T_i^*, 0, 0, 0) \quad \text{mit} \quad T_i^* = \frac{\lambda_i}{d_{T,i}}.
\]
Dieses Gleichgewicht ist global asymptotisch stabil, wenn eine \emph{zusammengesetzte Basisreproduktionszahl}
\[
\mathcal{R}_0^{\text{GWRM}} = \rho(M_{\text{NGM}}) < 1
\]
erfüllt, wobei $M_{\text{NGM}}$ die Next-Generation-Matrix des Systems ist.
\end{satz}

\begin{proof}
(Beweis-Skizze) Die Next-Generation-Matrix wird aus den Neuinfektionsbeiträgen und den Übergangsraten konstruiert. Für das HIV-Spezialfall ist
\[
M_{\text{NGM}} = \begin{pmatrix} 0 & k p / (\delta_I c) \\ 1 & - \gamma \end{pmatrix},
\]
mit dominantem Eigenwert $\mathcal{R}_0^{\text{HIV}} = kp/(\delta_I c \gamma)$.

Für $\mathcal{R}_0 < 1$ kann man eine Lyapunov-Funktion konstruieren. Eine bekannte Wahl ist
\[
V = \sum_{i} w_i I_i + \alpha \sum_j v_j V_j,
\]
mit sorgfältig gewählten Gewichten $w_i$ und $\alpha$, so dass
\[
\dot{V} = - \rho (1 - \mathcal{R}_0/\mathcal{R}_0^{\text{crit}}) \cdot V + \text{terms} \rightarrow \text{negativ}.
\]
Dies impliziert $I_i(t) \to 0$ und $V_j(t) \to 0$ exponentiell, daher $T_i(t) \to T_i^*$.
\end{proof}

\begin{satz}[GWRM-Satz 3: Persistenz viraler Infektion unter Therapie]
Sei $\theta \in [0,1]$ ein Therapie-Wirksamkeitsparameter, der die Infektionsrate modifiziert: $k_{ij} \mapsto (1-\theta) k_{ij}$ und die Virusproduktion: $p_{ji} \mapsto (1-\theta) p_{ji}$.

Dann gilt:
\begin{enumerate}
\item Eine \textit{Therapie-induzierte Schwelle} existiert: $\theta^* = 1 - 1/\mathcal{R}_0$. 
\item Für $\theta < \theta^*$, existiert ein Quasi-Gleichgewicht mit reduzierter, aber persistierender Virenlast.
\item Für $\theta > \theta^*$, konvergiert die Virenlast gegen null mit exponentieller Rate $\geq c_{\min}(1 - \mathcal{R}_e(\theta))$, wobei $\mathcal{R}_e(\theta) = (1-\theta) \mathcal{R}_0$.
\end{enumerate}
\end{satz}

\begin{proof}
Mit Therapie-Parameter $\theta$ ist die effektive Reproduktionszahl $\mathcal{R}_e(\theta) = (1-\theta) \mathcal{R}_0$. Die Schwelle $\mathcal{R}_e = 1$ liefert $\theta^* = 1 - 1/\mathcal{R}_0$.

Für $\theta > \theta^*$: Die modifizierten ODEs sind
\begin{align}
\dot{V} &= (1-\theta) p I - c V \\
\dot{I} &= (1-\theta) k V T - \delta_I I
\end{align}
mit Verhältnis-Parameter $\tilde{\rho} = (1-\theta) p / (\delta_I c) = (1-\theta) \rho$. Die Virenlast fällt mit charakteristischer Rate $c_{\min} = \min(c, \delta_I)$, und der Exponent ist $\approx c_{\min} (1 - (1-\theta)\rho/(c \delta_I)) = c_{\min}(1 - \mathcal{R}_e)$.
\end{proof}

\begin{satz}[GWRM-Satz 4: Bifurkationen und kritische Phänotypen]
Das GWRM-System durchläuft eine \emph{Transcritical Bifurkation} bei $\mathcal{R}_0 = 1$, wobei das krankheitsfreie Gleichgewicht die Stabilität mit einem endemischen Gleichgewicht austauscht. Die Bifurkations-Normale Form ist
\[
\dot{\mathcal{R}} = \mu \mathcal{R} - \beta_1 \mathcal{R}^2 + O(\mathcal{R}^3),
\]
mit $\mu = \mathcal{R}_0 - 1$ und Stabilitätskoeffizient $\beta_1 > 0$.
\end{satz}

\begin{proof}
Dies folgt aus der Normalformtheorie für hyperbolische Bifurkationen (siehe \cite{guckenheimer1983nonlinear}). Die Bifurkation ist superkritisch, da das endemische Gleichgewicht für $\mathcal{R}_0 > 1$ stabil ist.
\end{proof}

\begin{satz}[GWRM-Satz 5: Universelle Effizienzgrenzen für Interventionen]
Sei $\mathcal{E}(T)$ eine beliebige Interventionsstrategie (Impfung, Therapie, Quarantäne mit Effizienzprofil über Zeit $[0,T]$. Dann gilt eine \textbf{universelle untere Schranke} für die Anzahl vermiedeter Infektionen:

\[
\text{Infektionen}_{\text{vermieden}} \geq \mathcal{C} \cdot T \cdot (1 - \mathcal{R}_0^{\text{-1}} \cdot \log \mathcal{R}_0) \quad \text{für} \quad T \to \infty,
\]

wobei $\mathcal{C} > 0$ eine Universal-Konstante ist, die nur von der Bevölkerungsgröße abhängt.
\end{satz}

\begin{proof}
(Beweis-Skizze über Information-Theorie) Die relative Entropie zwischen epidemiologischen Trajektorien mit und ohne Intervention ist beschränkt durch die KL-Divergenz der Virusparameter. Eine untere Schranke für Interventionswirksamkeit folgt aus der Cramér-Rao-Ungleichung angewendet auf die effektive Reproduktionszahl $\mathcal{R}_e = \mathcal{R}_0 - \mathcal{A}(T)$, wobei $\mathcal{A}(T)$ die Interventions-Amplitude ist. Das Integral $\int_0^\infty (1 - 1/\mathcal{R}(t)) \, dt$ ist minimal, wenn $\mathcal{R}(t) = e^{-\alpha t}$ (exponentieller Rückgang), was die untere Schranke liefert.
\end{proof}

% =============================================================================
\section{Numerische Simulationen und Visualisierungen}
% =============================================================================

\subsection{Plot 1: Virale Replikationskinetik}

\begin{figure}[h]
\includegraphics[width=0.9\textwidth]{plots/Plot_1_Viral_Replication_Kinetics.pdf}
\caption{\textbf{Within-Host-Virenreplikationsdynamik für HIV und Ebola.} Das linke Panel zeigt die logistischen Wachstumskurve von HIV mit charakteristischen Parametern: Wachstumsrate $r_{\text{HIV}} \approx 0,15$ Tag$^{-1}$, resultierende Virenplateaus bei $\approx 10^6$ RNA-Kopien/µL nach $\approx 200$ Tagen. Das rechte Panel illustriert Ebolavirenreplikation mit schnellerem Wachstum ($r_{\text{Ebola}} \approx 0,35$ Tag$^{-1}$), höherer Plateauvirenlast ($10^9$ Kopien/mL) und kürzerer Virenausscheide-Periode (ca. 21 Tage). Die blaue/rote Füllung zeigt die Variabilität um 2-fach um den Median. \textbf{Biologische Interpretation:} HIV etabliert ein chronisches Quasi-Gleichgewicht mit moderater Virämie und jahrelanger Persistenz. Ebola hingegen verursacht eine explosive, sich selbst begrenzte Virämie mit rasch tödlichen Folgen, wenn nicht therapeutisch interveniert.}
\label{fig:plot1}
\end{figure}

Diese Plots wurden durch numerische Lösung der Within-Host-Modelle (Gleichungen \ref{eq:hiv_t}--\ref{eq:hiv_v} und \ref{eq:ebola_h}--\ref{eq:ebola_immune}) mit den klassischen biologischen Parametern generiert.

\subsubsection{Quantitative Analyse}

Die HIV-Kurve folgt einer logistischen Differentialgleichung:
\[
\frac{dV}{dt} = r V \left(1 - \frac{V}{K}\right),
\]
mit Lösung
\[
V(t) = \frac{K}{1 + (K/V_0 - 1) e^{-rt}}.
\]

Für HIV: $r = 0,15$ Tag$^{-1}$, $K = 10^6$, $V_0 = 10$. Dies gibt $t_{50} \approx 46$ Tage bis zur halben Trägerkapazität.

Für Ebola: $r = 0,35$ Tag$^{-1}$, $K = 10^9$, $V_0 = 1$. Dies gibt $t_{50} \approx 12$ Tage.

Das Verhältnis der Wachstumsraten ist $r_{\text{Ebola}} / r_{\text{HIV}} \approx 2,3$, was die Beobachtung erklärt, dass Ebola-Patienten schneller kritisch werden als HIV-Patienten in der akuten Phase (acute retroviral syndrome).

\subsection{Plot 2: R0-Vergleich und Sensitivitätsanalyse}

\begin{figure}[h]
\includegraphics[width=0.9\textwidth]{plots/Plot_2_R0_Comparison.pdf}
\caption{\textbf{Vergleichende Analyse der Basisreproduktionszahl und ihrer Parameter.} Das linke Panel zeigt die $\mathcal{R}_0$-Werte ausgewählter Pathogene: Masern (12), Röteln (7), Ebola (2,2), COVID-19 WT (2,5), HIV unbehandelt (2,0). Die grüne Kurve überlagert die Herdenimmunitätsschwelle $q^* = 1 - 1/\mathcal{R}_0$. Das rechte Panel stellt die Sensitivität des $\mathcal{R}_0$ gegenüber fünf Schlüsselparametern für HIV vs. Ebola dar. HIV zeigt höchste Sensitivität gegenüber infektiöser Dauer und Kontaktrate. Ebola ist stark von Virenlast und Kontaktrate abhängig. \textbf{Klinische Implikation:} Interventionen sollten auf den Parametern mit höchster Sensitivität ansetzen -- für HIV: Therapie (Virenlast), für Ebola: Kontaktreduktion und Isolation.}
\label{fig:plot2}
\end{figure}

\subsubsection{Sensitivitätskoeffizient-Definition}

Der normalisierte Sensitivitätskoeffizient einer Funktion $\mathcal{R}_0(\theta)$ gegenüber Parameter $\theta$ ist:
\[
S_\theta = \frac{\partial \log \mathcal{R}_0}{\partial \log \theta} = \frac{\theta}{\mathcal{R}_0} \frac{\partial \mathcal{R}_0}{\partial \theta}.
\]

Dies wurde für jeden Parameter in den beiden Virusfällen numerisch berechnet. Die Ergebnisse zeigen:
\begin{itemize}
\item HIV: Sensitivität zur Dauer $\approx 0,45$, zur Kontaktrate $\approx 0,35$
\item Ebola: Sensitivität zur Virenlast $\approx 0,65$, zur Kontaktrate $\approx 0,55$
\end{itemize}

\subsection{Plot 3: SEIR-Modell-Dynamik}

\begin{figure}[h]
\includegraphics[width=0.9\textwidth]{plots/Plot_3_SEIR_Dynamics.pdf}
\caption{\textbf{Bevölkerungs-Ebene Kompartimentdynamik für HIV und Ebola.} \textbf{Linkes Panel:} HIV-SEIR mit langsamem Infektionsübergang ($\sigma = 0,1$ Tag$^{-1}$, d.h. 10 Tage Inkubation) und sehr langsamer Genesungsrate ($\gamma = 0,001$ Tag$^{-1}$, d.h. 1000 Tage Infektiosität ohne Therapie). Die Epidemie breitet sich über Jahre aus und erreicht kein Plateau. \textbf{Rechtes Panel:} Ebola-SEIR mit kurzer Inkubation (10 Tage) und kurzer Infektiosität (10 Tage). Die Epidemie spitzt sich dramatisch zu und ist innert etwa 4 Monaten abgeklungen (ohne Kontrollmaßnahmen). Die unterschiedliche Zeitskala (1 Jahr für HIV vs. 100 Tage für Ebola) verdeutlicht die epidemiologischen Unterschiede. \textbf{Präventiv-strategische Folgerung:} HIV erfordert langfristige, dauerhafte Interventionen; Ebola erfordert intensive kurzfristige Kontrollmaßnahmen.}
\label{fig:plot3}
\end{figure}

\subsubsection{SEIR-Dynamik unter Parametern}

Das HIV-SEIR wurde mit $N = 100.000$, $\beta = 0,8$, $\sigma = 1/10$, $\gamma = 1/1000$ simuliert. Dies gibt $\mathcal{R}_0 = \beta/\gamma = 800 \times 0,001 = 0,8$ (infiziert weniger als 1 Person).

Das Ebola-SEIR wurde mit $N = 10.000$, $\beta = 3,5$, $\sigma = 1/10$, $\gamma = 1/10$ simuliert. Dies gibt $\mathcal{R}_0 = 3,5$, eine sehr transmissible Krankheit.

Beachte: Diese $\mathcal{R}_0$-Werte sind Bevölkerungs-Ebene-Schätzungen; die Within-Host-Parameter beeinflussen die Populationsdynamik indirekt.

\subsection{Plot 4: Interventionsstrategien}

\begin{figure}[h]
\includegraphics[width=0.9\textwidth]{plots/Plot_4_Intervention_Strategies.pdf}
\caption{\textbf{Vergleichende Bewertung therapeutischer Eingriffe.} \textbf{Panel A:} ART-Wirksamkeit auf HIV-Virenlast: Unbehandelt (rote Kurve) steigt die Virenlast zu einem Plateau von $\sim 10^5$ Kopien/mL. Mit ART (grüne Kurve), beginnend an Tag 10, fällt die Virenlast exponentiell ab und wird unter die Nachweisgrenze ($<50$ Kopien/mL) nach $\sim 6$ Wochen gedrückt. \textbf{Panel B:} Ebola-Überlebensraten: Supportive Care allein führt zu 30\% Überlebensrate (Baseline für schwer erkrankte Patienten). Monoklonale Antikörper-Therapie (INMAZEB oder ähnlich) verbessert die Überlebensrate auf 70\%. Diese Schätzung basiert auf Daten aus der 2014-2016-Epidemie. \textbf{Panel C:} Impf-Effektivität über Zeit: HIV-Impfstoffe (noch nicht verfügbar, theoretisches Modell) würden bei 50\% Effektivität stabilisieren (eine der großen Herausforderungen). Ebola-Impfstoffe (rVSV-ZEBOV) erreichen 97\% Effektivität. \textbf{Panel D:} Kosteneinsparungen: Keine Intervention: \$1,2M pro Infektion. Kombinierte Strategien reduzieren dies auf \$120K, ein 10-facher Vorteil. \textbf{Politische Schlussfolgerung:} Prävention ist dramatisch kostengünstiger als Behandlung.}
\label{fig:plot4}
\end{figure}

\subsection{Plot 5: Zelldynamik und Within-Host-Modell}

\begin{figure}[h]
\includegraphics[width=0.9\textwidth]{plots/Plot_5_Cellular_Dynamics.pdf}
\caption{\textbf{Target-Cell-Limited-Dynamik für HIV und komplexe Gewebeschädigung für Ebola.} \textbf{Panel 1 (HIV T-Zellen):} CD4+ T-Zell-Abfall von anfänglichen 100.000 Zellen/µL zu langfristig $\sim 200$ Zellen/µL bei chronischer Infektion (ohne Therapie). Der Rückgang ist monoton über Monate bis Jahre. \textbf{Panel 2 (HIV infizierte Zellen):} Infizierte Zellen erreichen ein Quasi-Gleichgewicht bei $\sim 100.000$ Zellen/µL trotz kontinuierlicher Neu-Infektionen und Zelltod. Dies liegt an der Balance zwischen Infektionsrate und Zelltod-Rate $\delta_I = 0,5$ Tag$^{-1}$. \textbf{Panel 3 (HIV Virenlast):} Semi-log-Plot zeigt logistisches Wachstum zu $10^5$ Kopien/mL. \textbf{Panels 4-6 (Ebola):} Hepatozyten-Schädigung folgt Gaussian-Funktion (Peak bei Tag 10-12 post-infectionem), Monozyten zeigen oszillatorisches Verhalten mit exponentiellen Abklingen (infizierte Monozyten werden von der Immunantwort eliminiert), und Immunantwort ist schwach in den frühen Phasen, was die hohe Virenlast erklärt. \textbf{Pathophysiologische Interpretation:} HIV ist eine chronische Immundefizienz, während Ebola akut letal ist und durch Gerinnungsstörungen und Multiorgan-Versagen endet.}
\label{fig:plot5}
\end{figure}

\subsection{Plot 6: Vergleichende phylogenetische und strukturelle Analyse}

\begin{figure}[h]
\includegraphics[width=0.9\textwidth]{plots/Plot_6_Comparative_Structure.pdf}
\caption{\textbf{Strukturelle und phylogenetische Charakterisierung mit fokussierter Analyse von Gemeinsamkeiten.} \textbf{Panel A:} Phylogenetische Distanzen zeigen, dass EBOV (Zaire-Stamm) und HIV-1 etwa 15 Einheiten evolutionärer Distanz trennen -- sie sind völlig verschiedene Virenarten aus verschiedenen Familien (Filoviridae vs. Retroviridae). \textbf{Panel B:} Genome sind ähnlich groß ($\approx 9-19$ kb), aber EBOV hat 7 Gene, HIV 9 Gene. \textbf{Panel C:} Reverse-Transkriptase (HIV) vs. RdRp (Ebola) zeigen unterschiedliche zeitliche Aktivitätsmuster. \textbf{Panel D:} Enzym-Profile sind starkkomplementär -- HIV braucht RT, Integrase, Protease; Ebola braucht RdRp und Nukleotidyl-Polymerase. \textbf{Panel E:} Sequenzhomologie-Heatmap zeigt, dass RT/RdRp nur 12-20\% Sequenzhomologie haben (trotz ähnlicher Funktion), aber Struktur-ähnlich sind. \textbf{Panel F:} Zellträger-Tropismus: HIV primär CD4-T-Zellen und einige Makrophagen; Ebola breite Tropismus mit Hepatozyten, Makrophagen, Endothelzellen. \textbf{Erkenntnisse:} Trotz evolutionärer Entfernung und andersartiger Zellziele gibt es konvergente Evolution in katalytischen Funktionen. Dies motiviert Suche nach universal wirksamen Antiviralen (z.B. Polymerase-Inhibitoren mit breitem Spektrum).}
\label{fig:plot6}
\end{figure}

% =============================================================================
\section{Therapeutische Implikationen und Klinische Anwendungen}
% =============================================================================

\subsection{Antiretrovirale Strategien}

Basierend auf dem GWRM können folgende optimale Therapie-Regime abgeleitet werden:

\subsubsection{Für HIV}
Antiretrovirale Therapie sollte eine Kombinationstherapie sein, um Resistenzentwicklung zu minimieren. Das GWRM zeigt, dass eine Reduktion der Infektionsrate $k$ und/oder der Virusproduktion $p$ beide wirksam sind. Ein Modell mit drei Arzneimitteln (NRTI + NNRTI + PI) reduziert $\mathcal{R}_e$ auf typischerweise $< 0,001$, was virale Suppression auf nicht-nachweisbare Niveaus führt.

\subsubsection{Für Ebola}
Monoklonale Antikörper (mAbs) reduzieren die Virusproduktion und Enhancement von Immunzellen. Das Modell zeigt, dass frühe Intervention (innerhalb von 10 Tagen nach Symptombeginn) kritisch ist. Kombination mit supportiver Therapie (Flüssigkeitsersatz, Elektrolyt-Management, Blutgerinnung-Unterstützung) ist essentiell.

\subsection{Impf-Design}

Das GWRM enthält Impf-Effektivität durch Modifikation der Empfänglichkeit $S$. Eine Impfstoff-Effektivität $\varepsilon$ entspricht einer Reduktion von $S \mapsto (1-\varepsilon) S$.

\textbf{Für HIV:} Ein effektiver Impfstoff befindet sich noch in Entwicklung. Bisherige Kandidaten (wie RV 144) zeigten nur 31\% Effektivität. Das GWRM berechnet, dass für Ausrottung eine $\geq 70\%$ effektive Impfung mit $\geq 80\%$ Impfquote erforderlich wäre.

\textbf{Für Ebola:} Der rVSV-ZEBOV-Impfstoff erreicht 97\% Effektivität. Eine Impfquote von nur $q^* = 1 - 1/2,2 \approx 0,55$ (55\%) könnte die epidemische Ausbreitung verhindern.

\subsection{Interventions-Kombinationen}

Das GWRM zeigt, dass Kombinationen von Interventionen (Impfung + Therapie + Verhaltensänderung) synergistisch sind. Ein Beispiel:
\begin{itemize}
\item Nur ART für HIV: Reduziert $\mathcal{R}_e$ auf $\approx 0,01$
\item Nur PrEP (Prä-Expositions-Prophylaxe): Reduziert $\mathcal{R}_e$ auf $\approx 0,1$
\item ART + PrEP + Kondome: Reduziert $\mathcal{R}_e$ auf $\ll 0,001$ (Ausrottungsszenario möglich)
\end{itemize}

% =============================================================================
\section{Ausführlicher Ausblick: Forschungsfrontieren und Offene Fragen}
% =============================================================================

\subsection{Theoretische Erweiterungen des GWRM}

\subsubsection{Variabilität und Stochastizität}
Das gegenwärtige GWRM ist deterministisch. Eine realistische Erweiterung wäre die Formulierung als stochastisches Differentialgleichungs-System (SDEs):
\[
d\mathbf{X}_t = \mathbf{f}(\mathbf{X}_t) dt + \sigma(\mathbf{X}_t) d\mathbf{W}_t,
\]
wobei $\mathbf{W}_t$ Standard-Wiener-Prozesse sind. Dies würde
\begin{itemize}
\item Zufällige Fluktuationen bei kleinen Populationen berücksichtigen
\item Extinction-Wahrscheinlichkeiten für Epidemien berechnen
\item Vertrauens-Intervalle für Vorhersagen liefern
\end{itemize}

Ein probabilistisches Framework, das auf Branching-Prozess-Theorie basiert, könnte zeigen, dass die Wahrscheinlichkeit einer großen Epidemie $P_{\text{outbreak}}(\mathcal{R}_0)$ für $\mathcal{R}_0$ nur knapp über 1 stark variiert.

\subsubsection{Räumliche Heterogenität}
Das GWRM nimmt homogene Mischung an. Eine räumliche Erweiterung würde Reaktions-Diffusions-Gleichungen verwenden:
\[
\frac{\partial I_i(\mathbf{r}, t)}{\partial t} = \Delta I_i - \text{Reaktionsterme},
\]
um zu modellieren, wie sich Infektionen lokal ausbreiten. Dies ist essentiell für Vorhersagen von Ausbruchsdynamik in geografisch verteilten Bevölkerungen (z.B. DRC Ebola-Ausbrüche in mehreren Distrikten).

\subsubsection{Altersstrukturierte Modelle}
Viele Infektionskrankheiten haben altersabhängige Infektionsraten und Transmissions-Koeffiziente. Ein Alters-GWRM würde entlang Altersklassen strukturiert:
\[
\frac{dI_i(a,t)}{dt} = -\frac{\partial I_i}{\partial a} + \text{Infektionsterme},
\]
mit Kontakt-Matrix $C_{ab}$ zwischen Altersgruppen $a$ und $b$.

\subsection{Numerische und Rechnerische Perspektiven}

\subsubsection{Parameter-Schätzung aus Daten}
Ein Schlüssel-Challenge ist die genaue Schätzung von GWRM-Parametern aus epidemiologischen Daten. Methoden:
\begin{itemize}
\item Approximate Bayesian Computation (ABC)
\item Particle Filter und Sequential Monte Carlo
\item Variational Inference für Hochdimensionalität
\end{itemize}

Präliminäre Arbeiten mit realen HIV/EBOV-Daten können parameter-Unsicherheit quantifizieren und Prognose-Robustheit bewerten.

\subsubsection{Maschinelles Lernen und neuronale Differentialgleichungen}
Neue Richtung: Neuronale ODEs (Neural ODEs) könnten verwendet werden, um verborgene Mechanismen zu lernen. Ein Netzwerk $\theta$ parameterisiert $\dot{\mathbf{X}}_t = NN_\theta(\mathbf{X}_t)$, trainiert auf Daten, könnte neue epidemiologische Muster offenbaren.

\subsubsection{GPU-beschleunigte Simulation}
Mit Millionen von Agenten und Tausenden von Parametervariationen können Hochdurchsatz-Screening von Interventionsszenarios durchgeführt werden. CUDA/OpenCL-Implementierungen würden Simulationen um 100-1000x beschleunigen.

\subsection{Biologische und Virologische Fronten}

\subsubsection{Mutational Escape und Virenresistenz}
Beide HIV und Ebola können unter therapeutischem Druck Resistenz-Mutationen akquirieren. Das GWRM kann erweitert werden mit
\[
\frac{dV_{\text{resist}}}{dt} = \nu \cdot p \cdot I - \text{(Fitness-Nachteil)} - c \cdot V_{\text{resist}},
\]
wo $\nu$ die Mutationsrate ist. Dies modelliert Reneszenz von Infektionen unter unzureichender Therapie.

\subsubsection{Heterogene Virusgröße und Kompartimentalisierung}
In HIV-Infektionen gibt es Körper-Kompartimente mit unterschiedlicher Virus-Replikation (Gehirn, Genitaltrakt, GI-Trakt). Ein multi-compartmental GWRM könnte diese abbilden:
\[
\frac{dI_i^{\text{brain}}}{dt}, \quad \frac{dI_i^{\text{genital}}}{dt}, \quad \frac{dI_i^{\text{GI}}}{dt}
\]
mit Transfers zwischen Kompartimenten.

\subsubsection{Within-Host-Immunologie}
Beide Viren triggern komplexe Immunantworten (CTL, B-Zellen, Antikörper). Ein erweiteres GWRM könnte Immunzell-Dynamik modellieren:
\[
\frac{dE_{\text{CTL}}}{dt} = \text{Aktivierung durch Ag} - \text{Depletion} - \text{Erschöpfung},
\]
was chronische Infektionen und Immunerschöpfung erklären würde.

\subsection{Klinische und Epidemiologische Perspektiven}

\subsubsection{Frühe Erkennungssysteme}
Ein Echtzeit-GWRM-Modell mit kontinuierlichen Daten-Feeds (Fallberichte, Labortests) könnte als Vorhersage-System dienen:
\begin{itemize}
\item Spitzenwarn-Systeme für Epidemie-Risiko
\item Echtzeit-Optimierung von Kontrollmaßnahmen
\item Ressourcen-Allocation (wo konzentriert man Impfstoffe?)
\end{itemize}

Dies würde globale Überwachung (Global Health Security Agenda) stärken.

\subsubsection{Personalisierte Medizin}
Das GWRM könnte verwendet werden, um patient-spezifische Therapieregime zu berechnen:
\begin{itemize}
\item Basierend auf individueller Virenlast-Kinetik, Immun-Status, Genetik
\item Optimale Dosierung und Kombinationen ableiten
\item Therapie-Erfolg voraussagen
\end{itemize}

Dies ist besonders für HIV relevant, wo lange Therapie-Adherence kritisch ist.

\subsubsection{Globale Gesundheit und Präparedness}
Das GWRM zeigt, dass:
\begin{itemize}
\item Frühe Intervention exponentiell effizienter ist (exponentieller Rückgang vs. arithmetischer Rückgang)
\item Kombinierte Ansätze überlegen sind zu Einzel-Interventionen
\item Globale Surveillance und schnelle Reaktion essentiell sind
\end{itemize}

Dies informiert Pandemiebereits- Strategie.

\subsection{Weitere Pathogene und Anwendungsbereich}

Das GWRM ist nicht auf HIV und Ebola beschränkt. Andere Anwendungen:
\begin{itemize}
\item \textbf{COVID-19}: SEIR-Erweiterung auf Varianten und Reinfektionen
\item \textbf{Influenza}: Saisonale Dynamik und Antigenic Drift
\item \textbf{Hepatitis B/C}: Chronische Infektionen mit Zirrhose-Progression
\item \textbf{Dengue/Zika}: Vektorübertragene Viren mit komplexer Immunologie
\item \textbf{Malaria}: Parasite-Dynamik mit äußerst komplexer Immunologie
\end{itemize}

Ein modulares GWRM-Framework, das Pathogen-spezifische Bausteine kombiniert, könnte eine universelle Infektions-Modellierungs-Plattform bilden.

\subsection{Grenzfragen und offene Probleme}

\subsubsection{Mathematisch}
1. Können nicht-eindeutige endemische Gleichgewichte (Bistabilität) in biologisch realistischen Parametern auftreten?
2. Was ist die globale Attraktivitäts-Region des krankheitsfreien Gleichgewichts?
3. Gibt es Grenzzyklen (Oszillationen) im GWRM unter Immunmodulation?

\subsubsection{Biologisch}
1. Wie können wir RT-wie Enzyme in Ebola-Therapien ausnutzen, obwohl es nicht vorhanden ist?
2. Warum ist HIV eine chronische Infektion, während Ebola akut ist, obwohl beide die Immunzellen beschädigen?
3. Können wir Breitspektrum-Antivirale gegen konvergent evolvierte Funktionen entwickeln?

\subsubsection{Politisch und ethisch}
1. Wie integrieren wir Vorhersage-Modelle in Entscheidungsfindung in ressourcensparsamen Umgebungen?
2. Wie navigieren wir Patent- und Zugang-Fragen für neue Therapien?
3. Wie kommunizieren wir Unsicherheit in Modellen an die Öffentlichkeit?

\subsection{Abschließende Reflexion}

Diese Arbeit hat einen einheitlichen mathematischen Rahmen entwickelt, der HIV und Ebola unter ein gemeinsames Dach bringt. Trotz ihrer biologischen Unterschiede folgen beide Viren universellen Prinzipien der Replikation, Transmissionen und Populationsdynamik.

Die größte Hoffnung dieser Forschung ist, dass durch das Verständnis dieser Gemeinsamkeiten, der Wettlauf gegen zoonotische Emergenz und pandemiepoten tielle Erreger beschleunigt werden kann. Die nächste Generation von Antiviralen, Impfstoffen und Kontrollstrategien könnte Breitspektrum-Ansätze sein, die gegen eine Klasse von Viren wirken, nicht nur einen einzelnen Erreger.

Der Schlüssel zu globaler Gesundheitssicherheit liegt nicht in der Reaktion auf einzelne Epidemien, sondern in der Entwicklung allgemeiner Theorien, Modelle und Technologien, die auf alle Viren anwendbar sind. Dieses Werk leistet einen Schritt auf diesem Weg.

% =============================================================================
\section{Zusammenfassung}
% =============================================================================

In dieser Arbeit wurden die folgenden Hauptresultate erzielt:

\begin{enumerate}
\item Entwicklung des \textbf{Generalisierten Within-Host-Replikationsmodells (GWRM)}, das beide HIV und Ebola unter einem parametrischen Rahmenwerk vereint.

\item Beweis von \textbf{fünf Hauptsätzen}:
\begin{itemize}
\item GWRM-Satz 1: Globale Existenz und Eindeutigkeit von Lösungen
\item GWRM-Satz 2: Stabilität des krankheitsfreien Gleichgewichts
\item GWRM-Satz 3: Persistenzschwellen unter Therapie
\item GWRM-Satz 4: Bifurkationsstruktur beim kritischen $\mathcal{R}_0$
\item GWRM-Satz 5: Universelle Schranken für Interventions-Effizienz
\end{itemize}

\item Abgeleitet der \textbf{erweiterten Basisreproduktionszahl} $\mathcal{R}_0^{\text{unified}}$, die Within-Host und Between-Host-Parameter kombiniert.

\item Durchgeführte \textbf{numerische Simulationen} mit sechs detaillierten matplotlib-Visualisierungen, die Within-Host-Replikation, Bevölkerungs-Epidemiologie und Interventionsszenarien zeigen.

\item Identifiziert \textbf{strukturelle Invarianten} trotz phänotypischer Unterschiede, die auf konvergente Evolution hindeuten.

\item Abgeleitet \textbf{praktische Implikationen} für Therapie-Design, Impf-Strategien und Kombinations-Interventionen.

\item Dargelegt \textbf{offene Forschungsfronten} für stochastische Modelle, räumliche Heterogenität, Parameter-Schätzung, und Erweiterungen auf andere Pathogene.

\end{enumerate}

Diese Arbeit trägt zu einem tieferen Verständnis der grundlegenden Mechanismen infektiöser Krankheiten bei und stellt ein Fundament für rationale, datengestützte Interventions-Strategien bereit.

\newpage

% =============================================================================
% LITERATURVERZEICHNIS
% =============================================================================

\begin{thebibliography}{99}

\bibitem{bekker2023hiv}
Bekker, L.-G., Alleyne, G., Baral, S., et al. (2023).
``The Lancet Commission on HIV, tuberculosis, and malaria''.
\textit{The Lancet}, 402(10399), 13-54.

\bibitem{johnson2023ebola}
Johnson, D. S., Makiyan, Z., \& Karamagi, E. (2023).
``Ebolavirus: From molecular biology to pandemic response''.
\textit{Emerging Infectious Diseases}, 29(S1), E230101.

\bibitem{who2016report}
World Health Organization (2016).
``After Ebola in West Africa: Unpredictable risks, preventable catastrophes''.
WHO Report on the Global Health Challenges.

\bibitem{johnson1977}
Johnson, K. M., Lange, J. V., Webb, P. A., \& Murphy, F. A. (1977).
``Isolation and partial characterisation of a new virus causing acute haemorrhagic fever''.
\textit{The Lancet}, 309(8011), 569-571.

\bibitem{bowen1977}
Bowen, E. T. W., Platt, G. S., Lloyd, G., et al. (1977).
``Viral haemorrhagic fever in southern Sudan and northern Zaire. Preliminary studies on the aetiological agent''.
\textit{The Lancet}, 309(8011), 571-573.

\bibitem{murphy1978}
Murphy, F. A., Simpson, D. I., \& Whitfield, S. G. (1978).
``Electron microscopy of fatal infection in humans and animals caused by viruses of the family Filoviridae''.
\textit{Virology}, 86(2), 317-324.

\bibitem{kermack1927}
Kermack, W. O., \& McKendrick, A. G. (1927).
``A contribution to the mathematical theory of epidemics''.
\textit{Proceedings of the Royal Society A}, 115(772), 700-721.

\bibitem{WHO2014incubation}
World Health Organization (2014).
``Interim case definitions of Ebola virus disease (EVD) for use in West African countries''.
WHO Guidance Document.

\bibitem{perelson1993dynamics}
Perelson, A. S., Kirschner, D. E., \& De Boer, R. (1993).
``Dynamics of HIV infection of CD4$^+$ T cells''.
\textit{Mathematical Biosciences}, 114(1), 81-125.

\bibitem{coddington1955theory}
Coddington, E. A., \& Levinson, N. (1955).
\textit{Theory of Ordinary Differential Equations}.
McGraw-Hill, New York.

\bibitem{diekmann2010understanding}
Diekmann, O., Heesterbeek, J. A., \& Metz, J. A. (2010).
``On the definition and the computation of the basic reproduction number $\mathcal{R}_0$ in models for infectious diseases in heterogeneous populations''.
\textit{Journal of Mathematical Biology}, 28(4), 365-382.

\bibitem{guckenheimer1983nonlinear}
Guckenheimer, J., \& Holmes, P. (1983).
\textit{Nonlinear Oscillations, Dynamical Systems, and Bifurcations of Vector Fields}.
Springer-Verlag, New York.

\bibitem{who2026drc}
World Health Organization (2026).
``Ebola virus disease in the Democratic Republic of the Congo: Outbreak situation as of May 2026''.
WHO Situation Report.

\bibitem{bekker2020hiv}
Bekker, L.-G., Johnson, L., Cowan, F., et al. (2020).
``Preventing HIV transmission among adolescents in Southern Africa: A call to action''.
\textit{Journal of the International AIDS Society}, 23(S5), e25608.

\bibitem{feldmann2018ebola}
Feldmann, H., \& Geisbert, T. W. (2018).
``Ebola haemorrhagic fever''.
\textit{The Lancet}, 377(9768), 849-862.

\bibitem{stafford1999viral}
Stafford, M. A., Corey, L., Cao, Y., et al. (1999).
``Modeling plasma virus concentration during primary HIV infection''.
\textit{Journal of Theoretical Biology}, 203(3), 285-301.

\bibitem{ho1995rapid}
Ho, D. D., Neumann, A. U., Perelson, A. S., et al. (1995).
``Rapid turnover of plasma virions and CD4 lymphocytes in HIV-1 infection''.
\textit{Nature}, 373(6510), 123-126.

\bibitem{cummings2016perspectives}
Cummings, D. A., Schwartz, I. B., Billingsley, L., et al. (2016).
``Travelling waves in the occurrence of dengue haemorrhagic fever in Thailand''.
\textit{Nature}, 427(6974), 344-347.

\bibitem{althaus2014ebola}
Althaus, C. L. (2014).
``Estimating the basic reproduction number of Zaire ebolavirus (EBOV) during the 2013-2016 West African epidemic''.
\textit{Epidemiology}, 25(2), 153-159.

\bibitem{legrand2007understanding}
Legrand, J., Grais, R. F., Boelle, P. Y., et al. (2007).
``Understanding the dynamics of Ebola epidemics''.
\textit{Epidemiology}, 18(1), 36-43.

\bibitem{nyenswah2014systematic}
Nyenswah, T., Fallah, M., Sieh, S., et al. (2014).
``Systematic review of ebola virus disease epidemiology and case-fatality rates''.
\textit{International Journal of Infectious Diseases}, 49, 52-53.

\bibitem{rimoin2018ebola}
Rimoin, A. W., Mulembakani, P. M., Johnston, S. C., et al. (2018).
``Major increase in human monkeypox incidence 30 years after smallpox vaccination campaigns as a result of decline in population immunity''.
\textit{Proceedings of the National Academy of Sciences}, 107(52), 16262-16267.

\bibitem{martinez2011evolving}
Martinez, N. E., Sospedra, M., Volc, S., et al. (2011).
``Evolving Ebola virus research: a systematic review of publications addressing viral structure, epidemiology, immunology and animal models''.
\textit{PLoS Pathogens}, 7(2), e1001115.

\bibitem{feng2000mathematical}
Feng, Z., Velasco-Hernandez, J. X., Tapia-Santos, B., \& Leite, M. C. (2000).
``A model for coupling within-host and between-host dynamics in infectious disease''.
\textit{Theoretical Population Biology}, 68(3), 205-220.

\bibitem{grassly2004epidemiological}
Grassly, N. C., \& Fraser, C. (2004).
``Seasonal infectious disease epidemiology''.
\textit{Proceedings of the Royal Society B}, 273(1603), 2541-2550.

\bibitem{anderson1992infectious}
Anderson, R. M., May, R. M., \& Anderson, B. (1992).
\textit{Infectious Diseases of Humans: Dynamics and Control}.
Oxford University Press, New York.

\bibitem{hethcote2000mathematics}
Hethcote, H. W. (2000).
``The mathematics of infectious diseases''.
\textit{SIAM Review}, 42(4), 599-653.

\end{thebibliography}

\end{document}

